% Generated draft: requires mathematical and source review.
\documentclass[11pt]{article}
\usepackage[margin=25mm]{geometry}
\usepackage{fontspec}
\setmainfont{Latin Modern Roman}
\usepackage{amsmath,amssymb,mathtools,unicode-math}
\setmathfont{Latin Modern Math}
\usepackage{enumitem,needspace,xurl,hyperref}
\setlength{\parindent}{0pt}
\setlength{\parskip}{5pt}
\setlength{\emergencystretch}{3em}
\newcommand{\N}{\mathbb{N}}
\newcommand{\Z}{\mathbb{Z}}
\newcommand{\Q}{\mathbb{Q}}
\newcommand{\R}{\mathbb{R}}
\newcommand{\C}{\mathbb{C}}
\newcommand{\F}{\mathbb{F}}
\newcommand{\A}{\mathbb{A}}
\newcommand{\PP}{\mathbb{P}}
\newcommand{\E}{\mathbb{E}}
\newcommand{\eps}{\varepsilon}
\newcommand{\dd}{\,\mathrm{d}}
\begin{document}
\section*{UnsolvedMath ID: 30005225}
\textbf{Source problem code:} catalogue entry
\subsection{Definitions and mathematical statement}
Let \(\Lambda_{\Q}\) be the ring of symmetric functions over \(\Q\), with Schur basis \(\{s_\rho\}\) indexed by partitions. For a partition \(\rho\), write \(|\rho|\) for its size, and let \(p_r\) denote the power-sum symmetric function of degree \(r\). For each integer \(r\geq 1\), let \(\psi_r\) be the \(\Q\)-algebra endomorphism of \(\Lambda_{\Q}\) determined by \(\psi_r(p_j)=p_{rj}\). Plethysm \(f[g]\) is obtained by expressing \(f\) as a polynomial in the power sums and substituting \(p_r\mapsto\psi_r(g)\).

For positive integers \(m,n\), partitions \(\lambda\vdash m\), \(\mu\vdash n\), and \(\nu\vdash mn\), use the convention
\[
 s_\lambda[s_\mu]
 =\sum_{\nu\vdash mn}a^\nu_{\lambda,\mu}s_\nu.
\]
These Schur expansion coefficients are the plethysm coefficients under this convention.

\textbf{Problem.} Give a general combinatorial description of \(a^\nu_{\lambda,\mu}\), uniformly for all positive integers \(m,n\) and all partitions \(\lambda\vdash m\), \(\mu\vdash n\), and \(\nu\vdash mn\). A customary interpretation of this request is to specify explicit finite combinatorial objects, with rules depending on \((\lambda,\mu,\nu)\), whose number equals \(a^\nu_{\lambda,\mu}\).

\textbf{Unresolved qualifications.} The supplied excerpt does not define the coefficient notation or specify the order of the two plethysm arguments; the displayed expansion makes the convention used in this draft explicit. It also does not define what qualifies as a combinatorial description, so the counting interpretation is an interpretation of the intended target. The statement above treats positive-degree Schur functions; the source does not specify whether empty partitions or degree-zero cases are intended.
\subsection{Short English statement}
Find a uniform combinatorial rule for the coefficients obtained when the plethysm of two Schur functions is expanded in the Schur basis. The draft makes the argument order and positive-degree domain explicit, but these conventions and the precise meaning of a combinatorial description require confirmation from the source.
\subsection{Sources}
\begin{itemize}
\item Local UnsolvedMath registry, supplied complete source record, permanent ID \(30005225\), source problem code \texttt{catalogue entry}, title \emph{Combinatorial Description of Plethysm Coefficients}. The supplied record labels the problem \texttt{open} and its statement \texttt{corrected\_verified}; these are supplied curation labels, not independent mathematical verification.
\item Supplied report citation: \emph{Character Theory and Categorification (2023)}, Oberwolfach Reports, DOI 10.4171/owr/2022/39, pp. 2267--2333. \url{https://doi.org/10.4171/owr/2022/39}. The record gives proposed year 2022 and also identifies the report as Oberwolfach Report 39/2022; these supplied year designations are retained without reconciliation.
\item The original excerpt identifies the general question as Problem 9 in Stanley's 1999 survey of positivity problems in algebraic combinatorics, cited there as reference [13]. The supplied record does not include the full bibliographic entry for [13]. The excerpt appears with the contribution title \emph{A recursive formula for plethysm coefficients and some applications}, attributed to Stacey Law, joint work with Y. Okitani. This does not establish that they originally proposed the general problem.
\item Supplied literature source: Paget--Wildon, \emph{two stability theorems}. \url{https://arxiv.org/abs/2505.06104}.
\item Supplied literature source: Bowman--Paget--Wildon, \emph{plethysm}. \url{https://arxiv.org/abs/2311.02721}.
\item Supplied literature source: Oberwolfach Report 39/2022. \url{https://ems.press/content/serial-article-files/46975}.
\item Literature assessment dated exactly \texttt{2026-08-22}: the supplied curation states that important formulas and special-case stability theorems exist, while the requested general combinatorial description remains open. This is supplied dated curation derived from the OWR record, not an independently verified literature assessment. The supplied background explicitly requests independent verification of time-sensitive status and cited claims.
\end{itemize}
\end{document}
